% Intended LaTeX compiler: pdflatex
\documentclass[a4paper,12pt]{article}
\usepackage[utf8]{inputenc}
\usepackage[T1]{fontenc}
\usepackage{graphicx}
\usepackage{longtable}
\usepackage{wrapfig}
\usepackage{rotating}
\usepackage[normalem]{ulem}
\usepackage{amsmath}
\usepackage{amssymb}
\usepackage{capt-of}
\usepackage{hyperref}
\usepackage[margin=0.2in]{geometry}
\usepackage{amsmath}
\usepackage{amsfonts}
\usepackage{indentfirst}
\usepackage{pgffor}
\usepackage{amssymb}
\usepackage{pifont}
\usepackage{cancel}
\usepackage{amsthm}
\usepackage{polynom}
\usepackage{tikz}
\usepackage[tikz]{bclogo}
\usepackage{listings}
\usepackage[most]{tcolorbox}
\usepackage{forest}
\usepackage{adjustbox}
\usepackage{tikz-3dplot}
\usepackage{mathtools}
\usepackage{pgfplots}
\usetikzlibrary{tikzmark,calc,fit,matrix,arrows,automata,positioning}
\usepackage{centernot}
\usepackage{lmodern}
\usepackage{physics}
\usepackage[boxed,linesnumbered,vlined]{algorithm2e}
\usepackage{algpseudocode}
\usepackage{algorithmicx}
\usepackage{svg}
\tikzstyle{every picture}+=[remember picture]
\DeclareMathOperator{\spn}{span}
\DeclareMathOperator{\dom}{domain}
\DeclareMathOperator{\ran}{range}
\DeclareMathOperator{\rng}{range}
\DeclareMathOperator{\img}{Im}
\DeclareMathOperator{\adj}{adj}
\DeclareMathOperator{\Var}{Var}
\DeclareMathOperator{\cov}{cov}
\newcommand\dif[1]{\:\textrm{d}#1}
\DeclarePairedDelimiter\ceil{\lceil}{\rceil}
\DeclarePairedDelimiter\floor{\lfloor}{\rfloor}
\newcommand{\ihat}{\hat{\textbf{\i}}}
\newcommand{\jhat}{\hat{\textbf{\j}}}
\newcommand{\khat}{\hat{\textbf{k}}}
\newcommand{\rhat}{\hat{\textbf{r}}}
\newcommand{\thetahat}{\boldsymbol{\hat{\theta}}}
\newcommand{\?}{\stackrel{?}{=}}
\newcommand{\cmark}{\text{\ding{51}}}
\newcommand{\xmark}{\text{\ding{55}}}
\author{mahmood}
\date{\today}
\title{probability tree}
\hypersetup{
 pdfauthor={mahmood},
 pdftitle={probability tree},
 pdfkeywords={},
 pdfsubject={},
 pdfcreator={Emacs 30.0.50 (Org mode 9.6)}, 
 pdflang={English}}
\begin{document}

\maketitle
\definecolor{Brown}{cmyk}{0,0.81,1,0.60}
\definecolor{OliveGreen}{cmyk}{0.64,0,0.95,0.40}
\definecolor{CadetBlue}{cmyk}{0.62,0.57,0.23,0}
\definecolor{lightlightgray}{gray}{0.9}
\lstset{
  language=C,                             % Code langugage
  basicstyle=\ttfamily,                   % Code font, Examples: \footnotesize, \ttfamily
  keywordstyle=\color{OliveGreen},        % Keywords font ('*' = uppercase)
  commentstyle=\color{gray},              % Comments font
  numbers=left,                           % Line nums position
  numberstyle=\tiny,                      % Line-numbers fonts
  stepnumber=1,                           % Step between two line-numbers
  numbersep=5pt,                          % How far are line-numbers from code
  backgroundcolor=\color{lightlightgray}, % Choose background color
  frame=single,                           % A frame around the code
  tabsize=2,                              % Default tab size
  captionpos=b,                           % Caption-position = bottom
  breaklines=true,                        % Automatic line breaking?
  breakatwhitespace=false,                % Automatic breaks only at whitespace?
  showspaces=false,                       % Dont make spaces visible
  showtabs=false,                         % Dont make tabls visible
  columns=flexible,                       % Column format
  morekeywords={__global__, __device__},  % CUDA specific keywords
  showstringspaces=false,                 % dont show spaces in strings
}
% so that latex doesnt complain about sage not being a language
\lstdefinelanguage{sage}{}

% environment definition for special blocks
\theoremstyle{definition}
\newtheorem{my_example}{Example}
\counterwithin{my_example}{section}
\counterwithin{my_example}{subsection}
\newtheorem{definition}{Definition}
\counterwithin{definition}{section}
\counterwithin{definition}{subsection}
\newtheorem{characteristic}{Characteristic}
\newtheorem{entailment}{Entailment}
%\newtheore*{proof}{Proof}
\newtheorem{lemma}{Lemma}
\newtheorem{question}{Question}
\newtheorem{subquestion}{Subquestion}[question]
%\newtheorem{note}{Note}
\newtheorem{answer}{Answer}
\counterwithin{answer}{question}
\counterwithin{answer}{subquestion}
\newtheorem{step}{Step}[answer]

% the following snippet auto resizes images to fit properly
% Determine if the image is too wide for the page.
% \makeatletter
% \def\ScaleIfNeeded{%
%   \ifdim\Gin@nat@width>\linewidth
%     \linewidth
%   \else
%     \Gin@nat@width
%   \fi
% }
% \makeatother
% % Resize figures that are too wide for the page.
% \let\oldincludegraphics\includegraphics
% \renewcommand\includegraphics[2][]{%
%   \oldincludegraphics[width=\ScaleIfNeeded]{#2}
% }

\newenvironment{note}
  {\begin{bclogo}[logo=\bcinfo, couleurBarre=orange, couleur=lightgray]{ note}}
  {\end{bclogo}}
\newenvironment{attention}
  {\begin{bclogo}[logo=\bcattention, couleurBarre=red, noborder=true, 
                couleur=red!15]{Important!}}
  {\end{bclogo}}

\begin{characteristic}
a tree consists of events, each event branches into the next events, a given path from event A to B would represent \(P(A \cap B)\) which is equal to \(P(A) \cdot P(B\mid A)\) and a log (or an arm, whatever) represents \href{probability.org}{conditional probability} \(P(B|A)\) (this is merely my understanding of it, dont take it for granted)\\[0pt]
\end{characteristic}
\begin{question}
in a given city all residents work either as fishers or in business, 70\% of the reidents are male and the rest are female, of the men 38\% work as fishers along with 10\% of women, we pick a person randomly\\[0pt]
\begin{subquestion}
what is the probability that the person is a male that works in fishing\\[0pt]
\begin{answer}
\begin{center}
\includesvg{/Users/mahmooz/brain/out/G1S0eQf}
\end{center}\\[0pt]
\end{answer}
\end{subquestion}
\begin{subquestion}
what is the probability that the person works in fishing\\[0pt]
\begin{answer}
\(P(\text{fisher}) = 0.7 \cdot 0.38 + 0.3 \cdot 0.1 = 0.296\)\\[0pt]
\end{answer}
\end{subquestion}
\begin{subquestion}
what is the probability that the person is a female who works in business\\[0pt]
\begin{answer}
\(P(\text{woman in business}) = 0.3 \cdot 0.9 = 0.27\)\\[0pt]
\end{answer}
\end{subquestion}
\begin{subquestion}
what is the probability that the person works in business\\[0pt]
\begin{answer}
\(P(\text{in business}) = 0.7 \cdot 0.62 + 0.3 \cdot 0.9 = 0.704\)\\[0pt]
\end{answer}
\end{subquestion}
\end{question}
\begin{question}
in a given population 40\% are men and the rest are women, of the men 10\% are unemployed, in total 13\% of the population is unemployed\\[0pt]
\begin{subquestion}
a person is chosen randomly, what is the probability that its a woman?\\[0pt]
\begin{answer}
\begin{center}
\includesvg{/Users/mahmooz/brain/out/rYPt0ig}
\end{center}\\[0pt]

\begin{center}
\includesvg{/Users/mahmooz/brain/out/mFjsNnA}
\end{center}\\[0pt]

\(0.13 = 0.4 \cdot 0.1 + 0.6p\) so:\\[0pt]

\begin{center}
\includesvg{/Users/mahmooz/brain/out/mFjsNnA}
\end{center}\\[0pt]

\[
P(\text{woman} \mid \text{unemployed}) = \frac{P(\text{unemployed woman})}{P(\text{unemployed})} = \frac{0.6 \cdot 0.15}{0.13} = 0.6923
\]\\[0pt]
\end{answer}
\end{subquestion}
\begin{subquestion}
we define the events A - an unemployed person is chosen, B - a man is chosen, are those events \href{probability.org}{independent}? are they \href{probability.org}{disjoint}?\\[0pt]
\begin{answer}
\[P(A \cap B) = P(\text{unemployed and man}) = 0.04 \neq 0\]\\[0pt]
A,B arent disjoint\\[0pt]
\[P(A) \cdot P(B) = 0.4 \cdot 0.13 = 0.052 \neq 0.04\]\\[0pt]
B,A are dependant\\[0pt]
\end{answer}
\end{subquestion}
\end{question}
\begin{question}
an employee asked her employer for a letter of recommendation for her new job, she knows if she gets a good recommendation letter she'll have a probability of 80\% to get the new job, 40\% if the letter is average and 10\% if the letter is bad, the probability that she gets a good letter is 0.7, an average letter 0.2 and a bad letter 0.1\\[0pt]
\begin{subquestion}
what is the probability that she'd get the job?\\[0pt]
\begin{answer}
\begin{center}
\includesvg{/Users/mahmooz/brain/out/JkUzci4}
\end{center}\\[0pt]

\begin{gather*}
  P(\text{accepted}) = 0.7 \cdot 0.8 + 0.2 \cdot 0.4 + 0.1 \cdot 0.1\\
  P(\text{good} \mid \text{accepted}) = \frac{P(\text{good and accepted})}{P(\text{accepted})} = \frac{0.7 \cdot 0.8}{0.65} = \frac{0.56}{0.65} = \frac{56}{65}
\end{gather*}
\end{answer}
\end{subquestion}
\begin{subquestion}
had she been accepted for the new job, what is the probability that she had received a good letter of recommendation?\\[0pt]
\end{subquestion}
\end{question}
\end{document}